\chapter{Chapter06 — Executable Finite Fields}
%%% generated by the blueprint pipeline from FABL.Chapter06.Constructions.ExecutableFiniteFields
\section{Overview}

Declarations of the Lean module \texttt{FABL.Chapter06.Constructions.ExecutableFiniteFields}.

\section{Declarations}

\begin{definition}[Enumeration of binary vectors]
\lean{FABL.allBinaryVectors}
\label{FABL.allBinaryVectors}
\leanok
\uses{FABL.F₂Cube}
Define a list of binary vectors recursively: in dimension zero it consists of the zero
vector; in dimension $\ell+1$, it consists first of the vectors obtained by prefixing
$0$ to each vector in the dimension-$\ell$ list, followed by those obtained by prefixing
$1$.
\end{definition}

\begin{theorem}[Elements of the two-element field]
\lean{FABL.f₂_eq_zero_or_one}
\leanok
\uses{FABL.𝔽₂}
\difficulty{4}
Every element $a$ of the field with two elements is either $0$ or $1$.
\end{theorem}

\begin{theorem}[Membership of every binary vector]
\lean{FABL.mem_allBinaryVectors}
\label{FABL.mem_allBinaryVectors}
\leanok
\uses{FABL.F₂Cube, FABL.allBinaryVectors, FABL.f₂_eq_zero_or_one}
\difficulty{5}
For every nonnegative integer $\ell$, every vector $x \in (\mathbb Z/2\mathbb Z)^\ell$
occurs in the list formed recursively from the unique zero-dimensional vector by
prefixing each vector at dimension $\ell$ by both possible bits to form the list at
dimension $\ell+1$.
\end{theorem}

\begin{theorem}[Distinct binary vectors]
\lean{FABL.allBinaryVectors_nodup}
\label{FABL.allBinaryVectors_nodup}
\leanok
\uses{FABL.allBinaryVectors, FABL.𝔽₂}
\difficulty{5}
For every natural number $\ell$, the recursively defined list of vectors in
$\mathbb{F}_2^\ell$ contains no repetitions.
\end{theorem}

\begin{theorem}[Number of binary vectors]
\lean{FABL.length_allBinaryVectors}
\label{FABL.length_allBinaryVectors}
\leanok
\uses{FABL.allBinaryVectors}
\difficulty{4}
For every nonnegative integer $\ell$, the list of all binary vectors of length $\ell$
over $\mathbb F_2$ has length $2^\ell$.
\end{theorem}

\begin{definition}[Enumeration of binary vectors]
\lean{FABL.binaryVectorEnumerationEquiv}
\label{FABL.binaryVectorEnumerationEquiv}
\leanok
\uses{FABL.F₂Cube, FABL.allBinaryVectors, FABL.allBinaryVectors_nodup, FABL.length_allBinaryVectors, FABL.mem_allBinaryVectors}
For each nonnegative integer $\ell$, the enumeration of binary vectors gives an
equivalence between the indices $0,\ldots,2^\ell-1$ and the vectors in $\{0,1\}^\ell$.
\end{definition}

\begin{theorem}[Value of the vector enumeration]
\lean{FABL.binaryVectorEnumerationEquiv_apply}
\label{FABL.binaryVectorEnumerationEquiv_apply}
\leanok
\uses{FABL.allBinaryVectors, FABL.binaryVectorEnumerationEquiv, FABL.length_allBinaryVectors}
\difficulty{1}
For every natural number $\ell$ and every index $i<2^\ell$, the binary vector
enumeration equivalence sends $i$ to the entry of the recursively defined list of all
binary vectors of length $\ell$ at the corresponding index under the canonical
identification of that list's length with $2^\ell$.
\end{theorem}

\begin{definition}[Enumeration of ordered vector pairs]
\lean{FABL.binaryVectorPairEnumerationEquiv}
\label{FABL.binaryVectorPairEnumerationEquiv}
\leanok
\uses{FABL.F₂Cube, FABL.binaryVectorEnumerationEquiv}
For each nonnegative integer $\ell$, there is an equivalence between the indices
$0,\ldots,2^\ell 2^\ell-1$ and the ordered pairs of length-$\ell$ binary vectors.
\end{definition}

\begin{definition}[Polynomial from binary coefficients]
\lean{FABL.binaryLowPolynomial}
\label{FABL.binaryLowPolynomial}
\leanok
\uses{FABL.F₂Cube, FABL.𝔽₂}
For a length-$\ell$ vector $a$ over $\mathbb{F}_2$, define the polynomial
$\sum_{i=0}^{\ell-1} a_i X^i$ over $\mathbb{F}_2$.
\end{definition}

\begin{definition}[Monic polynomial from binary coefficients]
\lean{FABL.binaryMonicPolynomial}
\label{FABL.binaryMonicPolynomial}
\leanok
\uses{FABL.F₂Cube, FABL.binaryLowPolynomial, FABL.𝔽₂}
For a length-$\ell$ vector $a$ over $\mathbb{F}_2$, define the polynomial
$X^\ell+\sum_{i=0}^{\ell-1}a_iX^i$ over $\mathbb{F}_2$.
\end{definition}

\begin{theorem}[Degree bound for the coefficient polynomial]
\lean{FABL.binaryLowPolynomial_degree_lt}
\label{FABL.binaryLowPolynomial_degree_lt}
\leanok
\uses{FABL.F₂Cube, FABL.binaryLowPolynomial}
\difficulty{1}
For every natural number $\ell$ and every coefficient vector $a\in\mathbb F_2^\ell$, the
degree of $\sum_{i=0}^{\ell-1}a_iX^i$ is strictly less than $\ell$.
\end{theorem}

\begin{theorem}[Monicity of a binary monic polynomial]
\lean{FABL.binaryMonicPolynomial_monic}
\label{FABL.binaryMonicPolynomial_monic}
\leanok
\uses{FABL.F₂Cube, FABL.binaryLowPolynomial_degree_lt, FABL.binaryMonicPolynomial}
\difficulty{1}
For every natural number $\ell$ and every coefficient vector $a\in\mathbb{F}_2^\ell$,
the polynomial $X^\ell+\sum_{i=0}^{\ell-1}a_iX^i$ is monic over $\mathbb{F}_2$.
\end{theorem}

\begin{theorem}[Degree of a binary monic polynomial]
\lean{FABL.binaryMonicPolynomial_degree}
\label{FABL.binaryMonicPolynomial_degree}
\leanok
\uses{FABL.F₂Cube, FABL.binaryLowPolynomial_degree_lt, FABL.binaryMonicPolynomial}
\difficulty{3}
For every natural number $\ell$ and every coefficient vector $a\in\mathbb{F}_2^\ell$,
the polynomial $X^\ell+\sum_{i=0}^{\ell-1}a_iX^i$ has degree $\ell$.
\end{theorem}

\begin{theorem}[Natural degree of a binary monic polynomial]
\lean{FABL.binaryMonicPolynomial_natDegree}
\label{FABL.binaryMonicPolynomial_natDegree}
\leanok
\uses{FABL.F₂Cube, FABL.binaryMonicPolynomial, FABL.binaryMonicPolynomial_degree}
\difficulty{1}
For every natural number $\ell$ and every coefficient vector $a\in\mathbb{F}_2^\ell$,
the polynomial $X^\ell+\sum_{i=0}^{\ell-1}a_iX^i$ has natural degree $\ell$.
\end{theorem}

\begin{theorem}[Polynomial from its coefficients]
\lean{FABL.binaryLowPolynomial_coefficients}
\label{FABL.binaryLowPolynomial_coefficients}
\leanok
\uses{FABL.binaryLowPolynomial, FABL.𝔽₂}
\difficulty{3}
Let $\ell$ be a natural number and let $p\in\mathbb F_2[X]$ have natural degree less
than $\ell$. The polynomial formed from its coefficients in degrees $0$ through $\ell-1$
is equal to $p$.
\end{theorem}

\begin{theorem}[Units in the binary polynomial ring]
\lean{FABL.isUnit_f₂Polynomial_iff_eq_one}
\leanok
\uses{FABL.𝔽₂}
\difficulty{4}
A polynomial over the field with two elements is a unit if and only if it is the
constant polynomial $1$.
\end{theorem}

\begin{definition}[First coordinate vector]
\lean{FABL.binaryOneCoefficients}
\label{FABL.binaryOneCoefficients}
\leanok
\uses{FABL.F₂Cube}
For each nonnegative integer $\ell$, define a length-$\ell$ vector over $\mathbb{F}_2$
whose coordinate at index $0$ is $1$ and whose other coordinates are $0$.
\end{definition}

\begin{definition}[Binary convolution coefficient]
\lean{FABL.binaryConvolutionCoefficient}
\label{FABL.binaryConvolutionCoefficient}
\leanok
\uses{FABL.F₂Cube, FABL.𝔽₂}
For a natural number $\ell$, let $a,b\in(\mathbb{Z}/2\mathbb{Z})^\ell$ and let
$d\in\mathbb{N}$. The binary convolution coefficient of degree $d$ is the sum in
$\mathbb{Z}/2\mathbb{Z}$ of $a_i b_j$ over pairs of natural numbers with $i+j=d$ and
$i,j<\ell$.
\end{definition}

\begin{definition}[Binary monic coefficient]
\lean{FABL.binaryMonicCoefficient}
\label{FABL.binaryMonicCoefficient}
\leanok
\uses{FABL.F₂Cube, FABL.𝔽₂}
For a natural number $\ell$, let $m\in(\mathbb{Z}/2\mathbb{Z})^\ell$ and
$d\in\mathbb{N}$. Define its monic coefficient at degree $d$ to be $m_d$ if $d<\ell$,
$1$ if $d=\ell$, and $0$ if $d>\ell$.
\end{definition}

\begin{definition}[Binary product coefficient equality]
\lean{FABL.binaryProductEqMonic}
\label{FABL.binaryProductEqMonic}
\leanok
\uses{FABL.F₂Cube, FABL.binaryConvolutionCoefficient, FABL.binaryMonicCoefficient}
For a natural number $\ell$, let $m,a,b\in(\mathbb{Z}/2\mathbb{Z})^\ell$. The Boolean
test returns true exactly when, for every integer $d$ with $0\le d\le 2\ell$, the
degree-$d$ binary convolution coefficient of $a$ and $b$ equals the degree-$d$ monic
coefficient specified by $m$.
\end{definition}

\begin{theorem}[Coefficients of the binary low polynomial]
\lean{FABL.coeff_binaryLowPolynomial}
\label{FABL.coeff_binaryLowPolynomial}
\leanok
\uses{FABL.F₂Cube, FABL.binaryLowPolynomial}
\difficulty{5}
Let $\ell$ be a nonnegative integer, let $a=(a_0,\ldots,a_{\ell-1})\in\mathbb F_2^\ell$,
and let $d$ be a nonnegative integer. The coefficient of $X^d$ in
$\sum_{i=0}^{\ell-1}a_iX^i$ is $a_d$ when $d<\ell$, and is $0$ when $d\ge\ell$.
\end{theorem}

\begin{theorem}[Constant polynomial from the first coordinate]
\lean{FABL.binaryLowPolynomial_binaryOneCoefficients}
\label{FABL.binaryLowPolynomial_binaryOneCoefficients}
\leanok
\uses{FABL.binaryLowPolynomial, FABL.binaryOneCoefficients, FABL.coeff_binaryLowPolynomial}
\difficulty{4}
Let $\ell$ be a positive integer, and let $a\in\mathbb F_2^\ell$ have value $1$ at
coordinate $0$ and value $0$ at every other coordinate. The polynomial
$\sum_{i=0}^{\ell-1} a_iX^i$ is the constant polynomial $1$.
\end{theorem}

\begin{theorem}[Characterization of the constant-one polynomial]
\lean{FABL.binaryLowPolynomial_eq_one_iff}
\label{FABL.binaryLowPolynomial_eq_one_iff}
\leanok
\uses{FABL.F₂Cube, FABL.binaryLowPolynomial, FABL.binaryLowPolynomial_binaryOneCoefficients, FABL.binaryOneCoefficients, FABL.coeff_binaryLowPolynomial, FABL.𝔽₂}
\difficulty{4}
Let $\ell$ be a positive integer and let $a\in\mathbb F_2^\ell$. The polynomial
$\sum_{i=0}^{\ell-1}a_iX^i$ equals $1$ if and only if $a_0=1$ and $a_i=0$ for every
$i>0$.
\end{theorem}

\begin{theorem}[Coefficient of a binary polynomial product]
\lean{FABL.binaryConvolutionCoefficient_eq_coeff_mul}
\label{FABL.binaryConvolutionCoefficient_eq_coeff_mul}
\leanok
\uses{FABL.F₂Cube, FABL.binaryConvolutionCoefficient, FABL.binaryLowPolynomial, FABL.coeff_binaryLowPolynomial}
\difficulty{4}
Let $a,b\colon\{0,1\}^{\ell}\to\mathbb{F}_2$ be binary vectors, and let
$d\in\mathbb{N}$. The coefficient of degree $d$ in the product of the polynomials whose
coefficients are the entries of $a$ and $b$ is the sum of $a_i b_j$ over all pairs of
indices with $i+j=d$, taking a term to be zero if either index is outside
$\{0,\ldots,\ell-1\}$.
\end{theorem}

\begin{theorem}[Polynomial Coefficient Formula]
\lean{FABL.binaryMonicCoefficient_eq_coeff}
\label{FABL.binaryMonicCoefficient_eq_coeff}
\leanok
\uses{FABL.F₂Cube, FABL.binaryMonicCoefficient, FABL.binaryMonicPolynomial, FABL.coeff_binaryLowPolynomial}
\difficulty{4}
Let $\ell$ be a natural number, let $m$ be a vector of length $\ell$ over
$\mathbb{F}_2$, and let $d$ be a natural number. The coefficient of degree $d$ in the
monic polynomial associated with $m$ equals the specified binary monic coefficient at
degree $d$.
\end{theorem}

\begin{theorem}[Polynomial product criterion]
\lean{FABL.binaryProductEqMonic_eq_true_iff}
\label{FABL.binaryProductEqMonic_eq_true_iff}
\leanok
\uses{FABL.F₂Cube, FABL.binaryConvolutionCoefficient, FABL.binaryConvolutionCoefficient_eq_coeff_mul, FABL.binaryLowPolynomial, FABL.binaryLowPolynomial_degree_lt, FABL.binaryMonicCoefficient, FABL.binaryMonicCoefficient_eq_coeff, FABL.binaryMonicPolynomial, FABL.binaryMonicPolynomial_degree, FABL.binaryProductEqMonic}
\difficulty{6}
Let $\ell$ be a positive integer, and let $m,a,b$ be binary vectors of length $\ell$.
The coefficient equations saying that the product of the two low-degree polynomials
associated with $a$ and $b$ has the coefficients specified by the monic polynomial
associated with $m$ hold if and only if those two polynomials are equal.
\end{theorem}

\begin{definition}[Proper binary factor witness]
\lean{FABL.binaryProperFactorWitness}
\label{FABL.binaryProperFactorWitness}
\leanok
\uses{FABL.F₂Cube, FABL.binaryOneCoefficients, FABL.binaryProductEqMonic}
For a natural number $\ell$, let $m,a,b\in(\mathbb{Z}/2\mathbb{Z})^\ell$. The Boolean
test returns true exactly when both $a$ and $b$ differ from the all-one coefficient
vector of length $\ell$, and their binary convolution coefficients agree with the monic
coefficients specified by $m$ in every degree from $0$ through $2\ell$.
\end{definition}

\begin{theorem}[Proper factor witness characterization]
\lean{FABL.binaryProperFactorWitness_eq_true_iff}
\label{FABL.binaryProperFactorWitness_eq_true_iff}
\leanok
\uses{FABL.F₂Cube, FABL.binaryLowPolynomial, FABL.binaryLowPolynomial_eq_one_iff, FABL.binaryMonicPolynomial, FABL.binaryProductEqMonic_eq_true_iff, FABL.binaryProperFactorWitness}
\difficulty{1}
For a positive integer $\ell$, let $m,a,b\in(\mathbb{F}_2)^\ell$, and let $A$ and $B$ be
the polynomials whose coefficients are given by $a$ and $b$, respectively, in degrees
$0$ through $\ell-1$. Let $M$ be the monic polynomial of degree $\ell$ whose
coefficients below degree $\ell$ are given by $m$. The binary proper factor witness for
$m,a,b$ equals true if and only if $A\ne 1$, $B\ne 1$, and $AB=M$.
\end{theorem}

\begin{definition}[Existence of a proper binary factor]
\lean{FABL.binaryHasProperFactor}
\label{FABL.binaryHasProperFactor}
\leanok
\uses{FABL.F₂Cube, FABL.allBinaryVectors, FABL.binaryProperFactorWitness}
For a natural number $\ell$ and a vector $m\in(\mathbb{Z}/2\mathbb{Z})^\ell$, the
Boolean test returns true exactly when there are vectors
$a,b\in(\mathbb{Z}/2\mathbb{Z})^\ell$, each different from the all-one coefficient
vector, whose binary convolution coefficients agree with the monic coefficients
specified by $m$ in every degree from $0$ through $2\ell$.
\end{definition}

\begin{theorem}[Characterization of proper polynomial factors]
\lean{FABL.binaryHasProperFactor_eq_true_iff}
\label{FABL.binaryHasProperFactor_eq_true_iff}
\leanok
\uses{FABL.F₂Cube, FABL.binaryHasProperFactor, FABL.binaryLowPolynomial, FABL.binaryMonicPolynomial, FABL.binaryProperFactorWitness_eq_true_iff, FABL.mem_allBinaryVectors}
\difficulty{1}
Let $\ell$ be a positive integer, and let $m$ be a vector of coefficients in $\mathbb{F}_2$ of
length $\ell$. The finite search for a proper factor of the monic polynomial determined
by $m$ returns true if and only if there are two vectors $a,b$ of coefficients in $\mathbb{F}_2$
of length $\ell$ such that their associated low-degree polynomials are each different
from $1$ and their product is the monic polynomial determined by $m$.
\end{theorem}

\begin{definition}[Binary monic irreducibility test]
\lean{FABL.binaryMonicIrreducibleTest}
\label{FABL.binaryMonicIrreducibleTest}
\leanok
\uses{FABL.F₂Cube, FABL.binaryHasProperFactor}
For a natural number $\ell$ and a vector $m\in(\mathbb{Z}/2\mathbb{Z})^\ell$, the
Boolean test returns true exactly when there are no vectors
$a,b\in(\mathbb{Z}/2\mathbb{Z})^\ell$, each different from the all-one coefficient
vector, whose binary convolution coefficients agree with the monic coefficients
specified by $m$ in every degree from $0$ through $2\ell$.
\end{definition}

\begin{theorem}[Irreducibility test for binary monic polynomials]
\lean{FABL.binaryMonicIrreducibleTest_eq_true_iff}
\label{FABL.binaryMonicIrreducibleTest_eq_true_iff}
\leanok
\uses{FABL.F₂Cube, FABL.binaryHasProperFactor, FABL.binaryHasProperFactor_eq_true_iff, FABL.binaryLowPolynomial, FABL.binaryLowPolynomial_coefficients, FABL.binaryMonicIrreducibleTest, FABL.binaryMonicPolynomial, FABL.binaryMonicPolynomial_monic, FABL.binaryMonicPolynomial_natDegree, FABL.isUnit_f₂Polynomial_iff_eq_one}
\difficulty{5}
Let $\ell$ be a positive integer and let $m\in(\mathbb{F}_2)^\ell$. The binary monic
irreducibility test on $m$ returns true if and only if the monic polynomial
$X^\ell+\sum_{i=0}^{\ell-1}m_iX^i$ is irreducible over $\mathbb{F}_2$.
\end{theorem}

\begin{definition}[Search for an irreducible binary modulus]
\lean{FABL.findIrreducibleBinaryModulus?}
\leanok
\uses{FABL.F₂Cube, FABL.allBinaryVectors, FABL.binaryMonicIrreducibleTest}
For a natural number $\ell$, search the vectors in $(\mathbb{Z}/2\mathbb{Z})^\ell$ for
one that passes the binary monic irreducibility test. Return the first such vector if
one exists, and return no vector otherwise.
\end{definition}

\begin{theorem}[Irreducible binary monic polynomial]
\lean{FABL.exists_irreducible_binaryMonicPolynomial}
\label{FABL.exists_irreducible_binaryMonicPolynomial}
\leanok
\uses{FABL.BinaryExtensionField, FABL.F₂Cube, FABL.binaryExtensionField_finrank, FABL.binaryLowPolynomial, FABL.binaryMonicPolynomial, FABL.𝔽₂}
\difficulty{4}
For every positive integer $\ell$, there is a polynomial over $\mathbb{F}_2$ of the
form $X^\ell+\sum_{i=0}^{\ell-1}a_iX^i$, with coefficients $a_i\in\mathbb{F}_2$,
that is irreducible.
\end{theorem}

\begin{theorem}[Existence of an irreducible binary modulus]
\lean{FABL.findIrreducibleBinaryModulus?_isSome}
\leanok
\uses{FABL.binaryMonicIrreducibleTest_eq_true_iff, FABL.exists_irreducible_binaryMonicPolynomial, FABL.findIrreducibleBinaryModulus?, FABL.mem_allBinaryVectors}
\difficulty{3}
For every positive natural number $\ell$, the search over binary coefficient vectors of
length $\ell$ returns a value; that is, its result is not empty.
\end{theorem}

\begin{definition}[Certified binary field implementation]
\lean{FABL.CertifiedBinaryFieldImplementation}
\label{FABL.CertifiedBinaryFieldImplementation}
\leanok
\uses{FABL.F₂Cube, FABL.binaryMonicPolynomial}
For a natural number $\ell$, a certified binary field implementation consists of a
vector $m\in(\mathbb{Z}/2\mathbb{Z})^\ell$ of low coefficients and a proof that the
monic polynomial specified by those coefficients is irreducible over
$\mathbb{Z}/2\mathbb{Z}$.
\end{definition}

\begin{definition}[Construction of an irreducible binary modulus]
\lean{FABL.buildCertifiedBinaryFieldImplementation}
\label{FABL.buildCertifiedBinaryFieldImplementation}
\leanok
\uses{FABL.CertifiedBinaryFieldImplementation, FABL.binaryMonicIrreducibleTest_eq_true_iff, FABL.findIrreducibleBinaryModulus?, FABL.findIrreducibleBinaryModulus?_isSome}
For each positive integer $\ell$, the construction assigns a coefficient vector
$m\in\mathbb{F}_2^\ell$ such that the monic polynomial
$X^\ell+\sum_{i=0}^{\ell-1}m_iX^i$ is irreducible over $\mathbb{F}_2$.
\end{definition}

\begin{definition}[Binary coefficient lookup]
\lean{FABL.binaryCoefficientAt}
\label{FABL.binaryCoefficientAt}
\leanok
\uses{FABL.F₂Cube, FABL.𝔽₂}
Given a coefficient vector $a\in\mathbb{F}_2^\ell$ and a nonnegative integer $d$, return
its $d$th coefficient when $d<\ell$, and return $0$ when $d\ge\ell$.
\end{definition}

\begin{definition}[Addition of binary coefficient vectors]
\lean{FABL.binaryAdd}
\label{FABL.binaryAdd}
\leanok
\uses{FABL.F₂Cube}
For coefficient vectors $a,b\in\mathbb{F}_2^\ell$, their sum is the vector whose $i$th
coefficient is $a_i+b_i$ in $\mathbb{F}_2$.
\end{definition}

\begin{definition}[Shifted binary coefficients]
\lean{FABL.binaryShiftUpCoefficient}
\label{FABL.binaryShiftUpCoefficient}
\leanok
\uses{FABL.F₂Cube, FABL.binaryCoefficientAt, FABL.𝔽₂}
Given a coefficient vector $a\in\mathbb{F}_2^\ell$ and a nonnegative integer $d$, the
shifted coefficient at position $d$ is $0$ if $d=0$, and otherwise is the coefficient of
$a$ at position $d-1$, with value $0$ when $d-1\ge\ell$.
\end{definition}

\begin{definition}[Top binary coefficient]
\lean{FABL.binaryTopCoefficient}
\label{FABL.binaryTopCoefficient}
\leanok
\uses{FABL.F₂Cube, FABL.𝔽₂}
For a positive integer $\ell$ and a coefficient vector $a\in\mathbb{F}_2^\ell$, the top
coefficient is $a_{\ell-1}$.
\end{definition}

\begin{definition}[Multiplication by $X$ modulo a binary polynomial]
\lean{FABL.binaryXMulMod}
\label{FABL.binaryXMulMod}
\leanok
\uses{FABL.F₂Cube, FABL.binaryAdd, FABL.binaryShiftUpCoefficient, FABL.binaryTopCoefficient}
Let $\ell$ be positive, and let $m,a\in\mathbb{F}_2^\ell$. Define a vector whose
coefficient at position $d$ is the shifted coefficient of $a$ at $d$ plus
$a_{\ell-1}m_d$, with addition and multiplication in $\mathbb{F}_2$.
\end{definition}

\begin{definition}[Scaled vector addition]
\lean{FABL.binaryScaleAdd}
\label{FABL.binaryScaleAdd}
\leanok
\uses{FABL.F₂Cube, FABL.binaryAdd, FABL.𝔽₂}
Given $b\in\mathbb{F}_2$ and vectors $p,q\in\mathbb{F}_2^\ell$, form the vector whose
$i$th coefficient is $q_i+bp_i$ in $\mathbb{F}_2$.
\end{definition}

\begin{definition}[Binary modular multiplication state]
\lean{FABL.binaryMulModState}
\label{FABL.binaryMulModState}
\leanok
\uses{FABL.F₂Cube, FABL.binaryCoefficientAt, FABL.binaryScaleAdd, FABL.binaryXMulMod}
Let $\ell$ be positive and let $m,x,y\in\mathbb{F}_2^\ell$. Define a sequence of pairs
of vectors by starting at $(0,x)$ and, for each $r\ge0$, replacing the current pair
$(u,v)$ by $(u+y_rv,\,Xv\bmod m)$, where $y_r=0$ for $r\ge\ell$ and multiplication by
$X$ modulo $m$ is the coefficient-vector operation that shifts $v$ upward and adds
$v_{\ell-1}m$.
\end{definition}

\begin{definition}[Binary multiplication modulo an irreducible polynomial]
\lean{FABL.binaryMulMod}
\label{FABL.binaryMulMod}
\leanok
\uses{FABL.CertifiedBinaryFieldImplementation, FABL.F₂Cube, FABL.binaryMulModState}
Let $\ell>0$, and let a certified binary field implementation of degree $\ell$ specify
an irreducible monic polynomial over $\mathbb{F}_2$ of the form
$X^\ell+\sum_{i=0}^{\ell-1}m_iX^i$. For binary vectors $x,y\in\mathbb{F}_2^\ell$, their
modular product is the vector returned as the first component of the multiplication
state after $\ell$ steps using that implementation.
\end{definition}

\begin{definition}[Binary vector encoding in an adjoining-root algebra]
\lean{FABL.binaryAdjoinRootEncode}
\label{FABL.binaryAdjoinRootEncode}
\leanok
\uses{FABL.F₂Cube, FABL.binaryLowPolynomial, FABL.binaryMonicPolynomial}
For binary vectors $m,a\in\mathbb{F}_2^\ell$, let
$f_m(X)=X^\ell+\sum_{i=0}^{\ell-1}m_iX^i$. The encoding of $a$ in the algebra obtained
by adjoining a root of $f_m$ is the class of $\sum_{i=0}^{\ell-1}a_iX^i$ modulo the
relation $f_m(X)=0$.
\end{definition}

\begin{theorem}[Coordinate formula for binary addition]
\lean{FABL.binaryAdd_apply}
\label{FABL.binaryAdd_apply}
\leanok
\uses{FABL.F₂Cube, FABL.binaryAdd}
\difficulty{1}
Let $\ell$ be a natural number, and let $a,b$ be vectors of length $\ell$ over
$\mathbb{F}_2$. At every coordinate $i$, their binary sum has value $a_i+b_i$.
\end{theorem}

\begin{theorem}[Coordinate formula for scale-add]
\lean{FABL.binaryScaleAdd_apply}
\label{FABL.binaryScaleAdd_apply}
\leanok
\uses{FABL.F₂Cube, FABL.binaryScaleAdd, FABL.𝔽₂}
\difficulty{1}
Let $\ell$ be a natural number, let $u\in\mathbb{F}_2$, and let
$p,c\in(\mathbb{F}_2)^\ell$. For every coordinate $i<\ell$, the $i$th coordinate of the
scale-add of $p$ by $u$ to $c$ is $c_i+u p_i$.
\end{theorem}

\begin{theorem}[Coefficients of the low-degree polynomial]
\lean{FABL.binaryCoefficientAt_eq_coeff_binaryLowPolynomial}
\label{FABL.binaryCoefficientAt_eq_coeff_binaryLowPolynomial}
\leanok
\uses{FABL.F₂Cube, FABL.binaryCoefficientAt, FABL.binaryLowPolynomial, FABL.coeff_binaryLowPolynomial}
\difficulty{1}
Let $\ell$ be a natural number, let $a$ be a binary vector of length $\ell$, and let $d$ be a
natural number. Form the polynomial $P(X)=\sum_{i=0}^{\ell-1}a_iX^i$ over the field of two
elements, and define $a_d=0$ for $d\ge\ell$. Then $a_d$ equals the coefficient of $X^d$ in
$P$.
\end{theorem}

\begin{theorem}[Coefficient of a shifted polynomial]
\lean{FABL.binaryShiftUpCoefficient_eq_coeff_X_mul}
\label{FABL.binaryShiftUpCoefficient_eq_coeff_X_mul}
\leanok
\uses{FABL.F₂Cube, FABL.binaryCoefficientAt_eq_coeff_binaryLowPolynomial, FABL.binaryLowPolynomial, FABL.binaryShiftUpCoefficient}
\difficulty{4}
Let $\ell$ be a natural number and let $a\in(\mathbb{F}_2)^\ell$. For every natural
number $d$, the coefficient of degree $d$ in $X A(X)$, where
$A(X)=\sum_{i=0}^{\ell-1}a_iX^i$, is $0$ when $d=0$ and otherwise is the coefficient of
$A(X)$ in degree $d-1$ (taken to be $0$ when $d-1\ge\ell$).
\end{theorem}

\begin{theorem}[Coefficient of multiplication by the indeterminate]
\lean{FABL.binaryCoefficientAt_binaryXMulMod}
\label{FABL.binaryCoefficientAt_binaryXMulMod}
\leanok
\uses{FABL.F₂Cube, FABL.binaryCoefficientAt, FABL.binaryMonicCoefficient, FABL.binaryShiftUpCoefficient, FABL.binaryTopCoefficient, FABL.binaryXMulMod}
\difficulty{5}
Let $\ell>0$, let $m$ and $a$ be binary vectors of length $\ell$, and let $d$ be a natural
number. Write $a_d$ for the $d$th coefficient of $a$, taking it to be zero when $d\ge
\ell$; write $a_{d-1}$ for the shifted coefficient, taking it to be zero when $d=0$; and
let $a_{\ell-1}$ be the top coefficient of $a$. Extend $m$ by setting its coefficient at
degree $\ell$ to $1$ and all coefficients above that to $0$. Then the $d$th coefficient of
the vector obtained by multiplying by the indeterminate modulo this monic polynomial is
$a_{d-1}+a_{\ell-1}m_d$, with arithmetic in the field of two elements.
\end{theorem}

\begin{theorem}[Polynomial of the shifted coefficients]
\lean{FABL.binaryLowPolynomial_binaryXMulMod}
\label{FABL.binaryLowPolynomial_binaryXMulMod}
\leanok
\uses{FABL.F₂Cube, FABL.binaryCoefficientAt_binaryXMulMod, FABL.binaryCoefficientAt_eq_coeff_binaryLowPolynomial, FABL.binaryLowPolynomial, FABL.binaryMonicCoefficient_eq_coeff, FABL.binaryMonicPolynomial, FABL.binaryShiftUpCoefficient_eq_coeff_X_mul, FABL.binaryTopCoefficient, FABL.binaryXMulMod}
\difficulty{3}
Let $\ell$ be a positive integer and let $m,a\in\mathbb F_2^\ell$. Form a new
coefficient vector by shifting the coefficients of $a$ up one degree, taking the
coefficient of degree $\ell-1$ in $a$, multiplying $m$ by that coefficient, and adding
the two vectors coordinatewise. Its associated polynomial is $X$ times the polynomial
associated to $a$, plus that top coefficient times $X^\ell$ plus the polynomial
associated to $m$.
\end{theorem}

\begin{theorem}[Additivity of the low polynomial]
\lean{FABL.binaryLowPolynomial_binaryAdd}
\label{FABL.binaryLowPolynomial_binaryAdd}
\leanok
\uses{FABL.F₂Cube, FABL.binaryAdd, FABL.binaryLowPolynomial, FABL.coeff_binaryLowPolynomial}
\difficulty{4}
For every dimension $\ell$ and every pair of binary coefficient vectors $a,b$ of length
$\ell$, the low-degree polynomial associated with their coordinatewise sum over
$\mathbb{F}_2$ is the sum of their associated low-degree polynomials.
\end{theorem}

\begin{theorem}[Polynomial of a scaled sum]
\lean{FABL.binaryLowPolynomial_binaryScaleAdd}
\label{FABL.binaryLowPolynomial_binaryScaleAdd}
\leanok
\uses{FABL.F₂Cube, FABL.binaryAdd, FABL.binaryLowPolynomial, FABL.binaryScaleAdd, FABL.coeff_binaryLowPolynomial, FABL.𝔽₂}
\difficulty{4}
Let $\ell$ be a natural number, let $b\in\mathbb F_2$, and let $p,a\in\mathbb F_2^\ell$.
The polynomial associated to the coordinatewise sum $a+bp$ is the polynomial associated
to $a$ plus $b$ times the polynomial associated to $p$.
\end{theorem}

\begin{theorem}[Encoding the zero vector]
\lean{FABL.binaryAdjoinRootEncode_zero}
\label{FABL.binaryAdjoinRootEncode_zero}
\leanok
\uses{FABL.F₂Cube, FABL.binaryAdjoinRootEncode, FABL.binaryLowPolynomial}
\difficulty{1}
Let $\ell$ be a nonnegative integer and $m$ a length-$\ell$ vector over $\mathbb F_2$.
The encoding of the zero vector in the quotient obtained by adjoining a root of
$X^\ell+\sum_{i<\ell}m_iX^i$ is zero.
\end{theorem}

\begin{theorem}[Addition under polynomial encoding]
\lean{FABL.binaryAdjoinRootEncode_binaryAdd}
\label{FABL.binaryAdjoinRootEncode_binaryAdd}
\leanok
\uses{FABL.F₂Cube, FABL.binaryAdd, FABL.binaryAdjoinRootEncode, FABL.binaryLowPolynomial_binaryAdd}
\difficulty{3}
Let $\ell$ be a natural number and let $m,a,b$ be vectors of length $\ell$ over
$\mathbb{F}_2$. Encoding the coordinatewise sum of $a$ and $b$ as a root of the monic
polynomial determined by $m$ gives the sum of their individual encodings in the
corresponding quotient algebra.
\end{theorem}

\begin{theorem}[Scale-and-add under polynomial encoding]
\lean{FABL.binaryAdjoinRootEncode_binaryScaleAdd}
\label{FABL.binaryAdjoinRootEncode_binaryScaleAdd}
\leanok
\uses{FABL.F₂Cube, FABL.binaryAdjoinRootEncode, FABL.binaryLowPolynomial_binaryScaleAdd, FABL.binaryMonicPolynomial, FABL.binaryScaleAdd, FABL.𝔽₂}
\difficulty{3}
Let $\ell$ be a natural number, let $m$ be a vector of length $\ell$ over
$\mathbb{F}_2$, and let $b$ be an element of $\mathbb{F}_2$. For vectors $u,v$ of length
$\ell$, encoding $v+bu$ using the monic polynomial determined by $m$ equals the encoding
of $v$ plus $b$ times the encoding of $u$ in the corresponding quotient algebra.
\end{theorem}

\begin{theorem}[Multiplication by the adjoined root]
\lean{FABL.binaryAdjoinRootEncode_binaryXMulMod}
\label{FABL.binaryAdjoinRootEncode_binaryXMulMod}
\leanok
\uses{FABL.F₂Cube, FABL.binaryAdjoinRootEncode, FABL.binaryLowPolynomial_binaryXMulMod, FABL.binaryMonicPolynomial, FABL.binaryXMulMod}
\difficulty{3}
Let $\ell>0$ and let $m,a$ be length-$\ell$ vectors over $\mathbb F_2$. Form the monic
polynomial $X^{\ell}+\sum_{i<\ell}m_iX^i$ and encode a vector by its low-coefficient
polynomial in the quotient obtained by adjoining a root of this monic polynomial. If
$a'$ is obtained by shifting the coefficients of $a$ up one degree, with constant
coefficient zero, and adding $a_{\ell-1}m$, then the encoding of $a'$ equals the
adjoined root times the encoding of $a$.
\end{theorem}

\begin{definition}[Prefix evaluation at an adjoined root]
\lean{FABL.binaryAdjoinRootPrefix}
\label{FABL.binaryAdjoinRootPrefix}
\leanok
\uses{FABL.F₂Cube, FABL.binaryCoefficientAt, FABL.binaryMonicPolynomial}
Let $m,y\in\mathbb{F}_2^\ell$, let $r$ be a natural number, and let $\alpha$ be a root
adjoined for $f_m(X)=X^\ell+\sum_{i=0}^{\ell-1}m_iX^i$. The prefix of length $r$ is the
sum $\sum_{i=0}^{r-1} c_i\alpha^i$ in the quotient algebra, where $c_i$ is the binary
coefficient of $y$ at index $i$.
\end{definition}

\begin{theorem}[Zero-length root-power prefix]
\lean{FABL.binaryAdjoinRootPrefix_zero}
\label{FABL.binaryAdjoinRootPrefix_zero}
\leanok
\uses{FABL.F₂Cube, FABL.binaryAdjoinRootPrefix}
\difficulty{1}
Let $\ell$ be a nonnegative integer, let $m,y$ be length-$\ell$ vectors over $\mathbb
F_2$, and let $\alpha$ be a root adjoined for $X^\ell+\sum_{i<\ell}m_iX^i$. The empty
sum of root powers with coefficients from $y$ is zero: $\sum_{i<0} y_i\alpha^i=0$.
\end{theorem}

\begin{theorem}[Successive root-power prefixes]
\lean{FABL.binaryAdjoinRootPrefix_succ}
\label{FABL.binaryAdjoinRootPrefix_succ}
\leanok
\uses{FABL.F₂Cube, FABL.binaryAdjoinRootPrefix, FABL.binaryCoefficientAt, FABL.binaryMonicPolynomial}
\difficulty{3}
Let $\ell$ be a nonnegative integer, let $m,y$ be length-$\ell$ vectors over $\mathbb
F_2$, and let $r$ be a nonnegative integer. In the quotient obtained by adjoining a root
$\alpha$ of $X^\ell+\sum_{i<\ell}m_iX^i$, the sum $\sum_{i<r+1} y_i\alpha^i$ equals
$\sum_{i<r} y_i\alpha^i+y_r\alpha^r$, where coefficients beyond index $\ell-1$ are taken
to be zero.
\end{theorem}

\begin{theorem}[Encoding as a root-power sum]
\lean{FABL.binaryAdjoinRootEncode_eq_prefix}
\label{FABL.binaryAdjoinRootEncode_eq_prefix}
\leanok
\uses{FABL.F₂Cube, FABL.binaryAdjoinRootEncode, FABL.binaryAdjoinRootPrefix, FABL.binaryCoefficientAt, FABL.binaryLowPolynomial, FABL.binaryMonicPolynomial}
\difficulty{4}
Let $\ell$ be a nonnegative integer, and let $m,y$ be length-$\ell$ vectors over
$\mathbb F_2$. In the quotient obtained by adjoining a root $\alpha$ of
$X^\ell+\sum_{i<\ell}m_iX^i$, the encoding of $y$ is $\sum_{i<\ell} y_i\alpha^i$.
\end{theorem}

\begin{theorem}[Correctness of the multiplication state]
\lean{FABL.binaryMulModState_correct}
\label{FABL.binaryMulModState_correct}
\leanok
\uses{FABL.F₂Cube, FABL.binaryAdjoinRootEncode, FABL.binaryAdjoinRootEncode_binaryScaleAdd, FABL.binaryAdjoinRootEncode_binaryXMulMod, FABL.binaryAdjoinRootPrefix, FABL.binaryAdjoinRootPrefix_succ, FABL.binaryMonicPolynomial, FABL.binaryMulModState}
\difficulty{4}
Let $\ell>0$, let $m,x,y\in\mathbb{F}_2^\ell$, and set
$f_m(X)=X^\ell+\sum_{i=0}^{\ell-1}m_iX^i$ and $\alpha=X$ in the quotient
$\mathbb{F}_2[X]/(f_m)$. For each vector $v$, write
$E(v)=\sum_{i=0}^{\ell-1}v_i\alpha^i$, and put $P_r(y)=\sum_{i=0}^{r-1}y_i\alpha^i$,
with $y_i=0$ for $i\ge\ell$. If the multiplication state at step $r$ is $(u_r,v_r)$,
then $E(u_r)=E(x)P_r(y)$ and $E(v_r)=\alpha^rE(x)$.
\end{theorem}

\begin{theorem}[Multiplication in the binary extension field]
\lean{FABL.adjoinRoot_mk_binaryMulMod}
\label{FABL.adjoinRoot_mk_binaryMulMod}
\leanok
\uses{FABL.CertifiedBinaryFieldImplementation, FABL.F₂Cube, FABL.binaryAdjoinRootEncode, FABL.binaryAdjoinRootEncode_eq_prefix, FABL.binaryLowPolynomial, FABL.binaryMonicPolynomial, FABL.binaryMulMod, FABL.binaryMulModState_correct}
\difficulty{3}
Let $\ell$ be a positive integer, and let $m,x,y\in(\mathbb{F}_2)^\ell$, where
$\mathbb{F}_2$ is the field with two elements. Suppose the monic polynomial
$X^\ell+\sum_{i=0}^{\ell-1}m_iX^i$ is irreducible over $\mathbb{F}_2$. The residue class
of the polynomial encoded by the modular product of $x$ and $y$ equals the product of
the residue classes of the polynomials encoded by $x$ and $y$ in the quotient by this
modulus.
\end{theorem}

\begin{definition}[Table of a binary-vector operation]
\lean{FABL.binaryOperationTable}
\label{FABL.binaryOperationTable}
\leanok
\uses{FABL.F₂Cube}
Given a list $E$ of binary vectors of length $\ell$ and a binary-vector operation
$\star$, form the list indexed by ordered pairs from $E\times E$, recording $(x,y,x\star
y)$ for each pair in the product-list order.
\end{definition}

\begin{definition}[Binary addition table]
\lean{FABL.binaryAdditionTable}
\label{FABL.binaryAdditionTable}
\leanok
\uses{FABL.F₂Cube, FABL.allBinaryVectors, FABL.binaryAdd, FABL.binaryOperationTable}
For a natural number $\ell$, the binary addition table lists $(x,y,x+y)$ for every
ordered pair of length-$\ell$ binary vectors, in the order induced by the list of all
such vectors.
\end{definition}

\begin{definition}[Binary multiplication table]
\lean{FABL.binaryMultiplicationTable}
\label{FABL.binaryMultiplicationTable}
\leanok
\uses{FABL.CertifiedBinaryFieldImplementation, FABL.F₂Cube, FABL.allBinaryVectors, FABL.binaryMulMod, FABL.binaryOperationTable}
Let $\ell>0$, and let a certified binary field implementation of degree $\ell$ specify
an irreducible monic binary polynomial for modular multiplication. The binary
multiplication table lists $(x,y,xy\bmod f)$ for every ordered pair of length-$\ell$
binary vectors, in the order induced by the list of all such vectors.
\end{definition}

\begin{definition}[Executable binary field model]
\lean{FABL.ExecutableBinaryFieldModel}
\label{FABL.ExecutableBinaryFieldModel}
\leanok
\uses{FABL.CertifiedBinaryFieldImplementation, FABL.F₂Cube}
An executable binary field model of degree $\ell$ consists of a certified binary field
implementation, a list of binary vectors representing the field elements, and ordered
addition and multiplication tables, each a list of input pairs together with their
outputs.
\end{definition}

\begin{definition}[Construction of an executable binary field model]
\lean{FABL.buildExecutableBinaryFieldModel}
\label{FABL.buildExecutableBinaryFieldModel}
\leanok
\uses{FABL.ExecutableBinaryFieldModel, FABL.allBinaryVectors, FABL.binaryAdditionTable, FABL.binaryMultiplicationTable, FABL.buildCertifiedBinaryFieldImplementation}
For $\ell>0$, construct an executable binary field model by choosing a certified
irreducible binary modulus of degree $\ell$, listing all length-$\ell$ binary vectors as
the field elements, and assembling their binary addition and modular multiplication
tables using that modulus.
\end{definition}

\begin{theorem}[Length of a binary operation table]
\lean{FABL.length_binaryOperationTable}
\label{FABL.length_binaryOperationTable}
\leanok
\uses{FABL.F₂Cube, FABL.binaryOperationTable}
\difficulty{1}
For any finite list $E$ of binary vectors of length $\ell$ over $\mathbb F_2$ and any
binary operation on such vectors, the list of triples $(a,b,a*b)$ indexed by ordered
pairs from $E$ has length $|E|^2$.
\end{theorem}

\begin{theorem}[Size of the binary addition table]
\lean{FABL.length_binaryAdditionTable}
\label{FABL.length_binaryAdditionTable}
\leanok
\uses{FABL.binaryAdditionTable, FABL.length_allBinaryVectors, FABL.length_binaryOperationTable}
\difficulty{3}
For every nonnegative integer $\ell$, the list of triples $(a,b,a+b)$ formed from all
ordered pairs of binary vectors of length $\ell$ over $\mathbb F_2$ has length
$2^{2\ell}$.
\end{theorem}

\begin{theorem}[Size of the binary multiplication table]
\lean{FABL.length_binaryMultiplicationTable}
\label{FABL.length_binaryMultiplicationTable}
\leanok
\uses{FABL.CertifiedBinaryFieldImplementation, FABL.binaryMultiplicationTable, FABL.length_allBinaryVectors, FABL.length_binaryOperationTable}
\difficulty{3}
Let $\ell>0$ and choose a monic irreducible polynomial of degree $\ell$ over $\mathbb
F_2$, represented by its $\ell$ lower coefficients. The list of triples $(a,b,ab)$ for
all ordered pairs of binary vectors of length $\ell$, with multiplication performed
modulo this polynomial by the specified implementation, has length $2^{2\ell}$.
\end{theorem}

\begin{theorem}[Membership in the binary addition table]
\lean{FABL.mem_binaryAdditionTable}
\label{FABL.mem_binaryAdditionTable}
\leanok
\uses{FABL.F₂Cube, FABL.binaryAdd, FABL.binaryAdditionTable, FABL.binaryOperationTable, FABL.mem_allBinaryVectors}
\difficulty{3}
For every nonnegative integer $\ell$ and every $x,y\in\mathbb F_2^\ell$, the triple
$(x,y,x+y)$ belongs to the table of all ordered pairs of binary vectors together with
their coordinatewise sum.
\end{theorem}

\begin{theorem}[Membership in the binary multiplication table]
\lean{FABL.mem_binaryMultiplicationTable}
\label{FABL.mem_binaryMultiplicationTable}
\leanok
\uses{FABL.CertifiedBinaryFieldImplementation, FABL.F₂Cube, FABL.binaryMulMod, FABL.binaryMultiplicationTable, FABL.binaryOperationTable, FABL.mem_allBinaryVectors}
\difficulty{3}
Let $\ell>0$, and let a certified binary field implementation of dimension $\ell$ be
given. For all binary vectors $x,y\in\mathbb{F}_2^\ell$, the triple consisting of $x$,
$y$, and their modular product belongs to the multiplication table formed from all
binary vectors of that dimension.
\end{theorem}

\begin{theorem}[Completeness of the binary field model]
\lean{FABL.buildExecutableBinaryFieldModel_complete}
\label{FABL.buildExecutableBinaryFieldModel_complete}
\leanok
\uses{FABL.F₂Cube, FABL.binaryAdd, FABL.binaryMulMod, FABL.buildCertifiedBinaryFieldImplementation, FABL.buildExecutableBinaryFieldModel, FABL.mem_allBinaryVectors, FABL.mem_binaryAdditionTable, FABL.mem_binaryMultiplicationTable}
\difficulty{1}
Let $\ell>0$, and let $x,y$ be binary vectors of length $\ell$ over the
field with two elements. In the constructed executable binary field model, $x$ occurs in
the element list, the triple consisting of $x$, $y$, and their coordinatewise sum occurs
in the addition table, and the triple consisting of $x$, $y$, and their modular product
using the model's certified modulus occurs in the multiplication table.
\end{theorem}

\begin{definition}[List traversal work]
\lean{FABL.binaryConstructorTraversalWork}
\label{FABL.binaryConstructorTraversalWork}
\leanok
For a type of elements $\alpha$ and a natural number $s$ giving the work per step,
define the traversal work of a list recursively: the empty list has work $0$, and a
nonempty list has work $s$ plus the traversal work of its tail.
\end{definition}

\begin{theorem}[Traversal work for a list]
\lean{FABL.binaryConstructorTraversalWork_eq}
\label{FABL.binaryConstructorTraversalWork_eq}
\leanok
\uses{FABL.binaryConstructorTraversalWork}
\difficulty{4}
For any type of items, any natural number $s$ representing the work per step, and any
finite list of items, the total traversal work is the list length multiplied by $s$.
\end{theorem}

\begin{definition}[Binary vector enumeration work]
\lean{FABL.binaryVectorEnumerationWork}
\label{FABL.binaryVectorEnumerationWork}
\leanok
\uses{FABL.allBinaryVectors}
Let $V_\ell$ be the list of binary vectors of length $\ell$, defined by $V_0=[0]$ and,
for each $\ell$, by concatenating the vectors obtained by prefixing $0$ to each element
of $V_\ell$ with those obtained by prefixing $1$. Define the enumeration work by $W_0=1$
and $W_{\ell+1}=W_\ell+3|V_\ell|+1$, where $|V_\ell|$ is the list length.
\end{definition}

\begin{definition}[Binary addition work]
\lean{FABL.binaryAddWork}
\label{FABL.binaryAddWork}
\leanok
\uses{FABL.binaryConstructorTraversalWork}
For each natural number $\ell$, define the binary addition work to be the traversal work
with unit cost per step of the list of natural numbers from $0$ through $\ell-1$.
\end{definition}

\begin{definition}[Binary multiplication reduction work]
\lean{FABL.binaryXMulModWork}
\label{FABL.binaryXMulModWork}
\leanok
\uses{FABL.binaryConstructorTraversalWork}
For each natural number $\ell$, define the binary multiplication reduction work to be
one plus the traversal work, at cost $4$ per step, of the list of natural numbers from
$0$ through $\ell-1$.
\end{definition}

\begin{definition}[Binary scalar addition work]
\lean{FABL.binaryScaleAddWork}
\label{FABL.binaryScaleAddWork}
\leanok
\uses{FABL.binaryConstructorTraversalWork}
For each natural number $\ell$, define the binary scalar addition work to be the
traversal work, at cost $2$ per step, of the list of natural numbers from $0$ through
$\ell-1$.
\end{definition}

\begin{definition}[Binary multiplication state work]
\lean{FABL.binaryMulModStateWork}
\label{FABL.binaryMulModStateWork}
\leanok
\uses{FABL.binaryScaleAddWork, FABL.binaryXMulModWork}
For natural numbers $\ell$ and $r$, define the state work recursively by $S_{\ell,0}=0$
and $S_{\ell,r+1}=S_{\ell,r}+A_\ell+M_\ell+1$, where $A_\ell$ is the binary scalar
addition work and $M_\ell$ is the binary multiplication reduction work.
\end{definition}

\begin{definition}[Binary multiplication work]
\lean{FABL.binaryMulModWork}
\label{FABL.binaryMulModWork}
\leanok
\uses{FABL.binaryMulModStateWork}
For each natural number $\ell$, define the binary multiplication work to be the binary
multiplication state work at state index $\ell$.
\end{definition}

\begin{definition}[Binary arithmetic work]
\lean{FABL.binaryArithmeticWork}
\label{FABL.binaryArithmeticWork}
\leanok
\uses{FABL.binaryAddWork, FABL.binaryMulModWork}
For each natural number $\ell$, define the binary arithmetic work to be the maximum of
the binary addition work and the binary multiplication work.
\end{definition}

\begin{theorem}[Binary addition work]
\lean{FABL.binaryAddWork_eq}
\label{FABL.binaryAddWork_eq}
\leanok
\uses{FABL.binaryAddWork, FABL.binaryConstructorTraversalWork_eq}
\difficulty{1}
For every natural number $\ell$, the work assigned to binary addition on vectors of
length $\ell$ is $\ell$.
\end{theorem}

\begin{theorem}[Binary multiplication work]
\lean{FABL.binaryXMulModWork_eq}
\label{FABL.binaryXMulModWork_eq}
\leanok
\uses{FABL.binaryConstructorTraversalWork_eq, FABL.binaryXMulModWork}
\difficulty{1}
For every natural number $\ell$, the work required by the binary multiplication
procedure is exactly $4\ell+1$.
\end{theorem}

\begin{theorem}[Scale-add work count]
\lean{FABL.binaryScaleAddWork_eq}
\label{FABL.binaryScaleAddWork_eq}
\leanok
\uses{FABL.binaryConstructorTraversalWork_eq, FABL.binaryScaleAddWork}
\difficulty{1}
For every natural number $\ell$, the scale-add work is $2\ell$.
\end{theorem}

\begin{theorem}[Work of the multiplication state]
\lean{FABL.binaryMulModStateWork_eq}
\label{FABL.binaryMulModStateWork_eq}
\leanok
\uses{FABL.binaryMulModStateWork, FABL.binaryScaleAddWork_eq, FABL.binaryXMulModWork_eq}
\difficulty{4}
For all natural numbers $\ell$ and $r$, the work assigned to the multiplication state
after $r$ steps is $r(6\ell+2)$.
\end{theorem}

\begin{theorem}[Work of binary modular multiplication]
\lean{FABL.binaryMulModWork_eq}
\label{FABL.binaryMulModWork_eq}
\leanok
\uses{FABL.binaryMulModStateWork_eq, FABL.binaryMulModWork}
\difficulty{1}
For every natural number $\ell$, the work assigned to binary modular multiplication is
$\ell(6\ell+2)$.
\end{theorem}

\begin{theorem}[Explicit arithmetic work bound]
\lean{FABL.binaryArithmeticWork_le}
\label{FABL.binaryArithmeticWork_le}
\leanok
\uses{FABL.binaryAddWork_eq, FABL.binaryArithmeticWork, FABL.binaryMulModWork_eq}
\difficulty{3}
For every natural number $\ell$, the binary arithmetic work is at most $8(\ell+1)^2$.
\end{theorem}

\begin{definition}[Binary product comparison work]
\lean{FABL.binaryProductComparisonWork}
\label{FABL.binaryProductComparisonWork}
\leanok
\uses{FABL.binaryConstructorTraversalWork}
For each natural number $\ell$, define the product comparison work to be the constructor
traversal work with budget $2\ell+2$ on the list of integers from $0$ through $2\ell$.
\end{definition}

\begin{definition}[Binary factor search work]
\lean{FABL.binaryFactorSearchWork}
\label{FABL.binaryFactorSearchWork}
\leanok
\uses{FABL.allBinaryVectors, FABL.binaryConstructorTraversalWork, FABL.binaryProductComparisonWork}
For each natural number $\ell$, define the factor search work to be the constructor
traversal work with budget equal to the product comparison work plus $4$, on the set of
ordered pairs of binary vectors of length $\ell$.
\end{definition}

\begin{definition}[Binary modulus search work]
\lean{FABL.binaryModulusSearchWork}
\label{FABL.binaryModulusSearchWork}
\leanok
\uses{FABL.allBinaryVectors, FABL.binaryConstructorTraversalWork, FABL.binaryFactorSearchWork}
For each natural number $\ell$, define the modulus search work to be the constructor
traversal work with budget equal to the factor search work plus $1$, on the set of
binary vectors of length $\ell$.
\end{definition}

\begin{definition}[Binary operation table work]
\lean{FABL.binaryOperationTableWork}
\label{FABL.binaryOperationTableWork}
\leanok
\uses{FABL.allBinaryVectors, FABL.binaryConstructorTraversalWork}
For a natural number $\ell$ and a natural number $w$, define the operation table work to
be the constructor traversal work with budget $w+1$, on the set of ordered pairs of
binary vectors of length $\ell$.
\end{definition}

\begin{definition}[Binary addition table work]
\lean{FABL.binaryAdditionTableWork}
\label{FABL.binaryAdditionTableWork}
\leanok
\uses{FABL.binaryAddWork, FABL.binaryOperationTableWork}
For each natural number $\ell$, define the addition table work to be the operation table
work with operation work given by the binary addition work at parameter $\ell$.
\end{definition}

\begin{definition}[Binary multiplication table work]
\lean{FABL.binaryMultiplicationTableWork}
\label{FABL.binaryMultiplicationTableWork}
\leanok
\uses{FABL.binaryMulModWork, FABL.binaryOperationTableWork}
For each natural number $\ell$, define the multiplication table work to be the operation
table work with operation work given by the binary multiplication modulo work at
parameter $\ell$.
\end{definition}

\begin{definition}[Binary field representation bits]
\lean{FABL.binaryFieldRepresentationBits}
\label{FABL.binaryFieldRepresentationBits}
\leanok
\uses{FABL.allBinaryVectors}
For each natural number $\ell$, define the field representation bit count to be $\ell$
times one plus the number of binary vectors of length $\ell$, plus six times the number
of ordered pairs of such vectors.
\end{definition}

\begin{definition}[Binary field preprocessing work]
\lean{FABL.binaryFieldPreprocessingWork}
\label{FABL.binaryFieldPreprocessingWork}
\leanok
\uses{FABL.binaryAdditionTableWork, FABL.binaryFieldRepresentationBits, FABL.binaryModulusSearchWork, FABL.binaryMultiplicationTableWork, FABL.binaryVectorEnumerationWork}
For each natural number $\ell$, define the field preprocessing work as the sum of the
binary vector enumeration work, modulus search work, addition table work, multiplication
table work, and field representation bit count, all at parameter $\ell$.
\end{definition}

\begin{theorem}[Number of ordered binary vector pairs]
\lean{FABL.length_binaryVectorPairs}
\label{FABL.length_binaryVectorPairs}
\leanok
\uses{FABL.allBinaryVectors, FABL.length_allBinaryVectors}
\difficulty{3}
For every nonnegative integer $\ell$, the list of ordered pairs of binary vectors of
length $\ell$ over $\mathbb F_2$ has length $2^{2\ell}$.
\end{theorem}

\begin{theorem}[Binary vector enumeration work bound]
\lean{FABL.binaryVectorEnumerationWork_le}
\label{FABL.binaryVectorEnumerationWork_le}
\leanok
\uses{FABL.binaryVectorEnumerationWork, FABL.length_allBinaryVectors}
\difficulty{4}
For every integer $\ell\ge 0$, the work required to enumerate all binary vectors of
length $\ell$ is at most $4(\ell+1)2^\ell$.
\end{theorem}

\begin{theorem}[Product comparison work formula]
\lean{FABL.binaryProductComparisonWork_eq}
\label{FABL.binaryProductComparisonWork_eq}
\leanok
\uses{FABL.binaryConstructorTraversalWork_eq, FABL.binaryProductComparisonWork}
\difficulty{1}
For every natural number $\ell$, the product comparison work is $(2\ell+1)(2\ell+2)$.
\end{theorem}

\begin{theorem}[Binary factor search work]
\lean{FABL.binaryFactorSearchWork_eq}
\label{FABL.binaryFactorSearchWork_eq}
\leanok
\uses{FABL.binaryConstructorTraversalWork_eq, FABL.binaryFactorSearchWork, FABL.binaryProductComparisonWork, FABL.length_binaryVectorPairs}
\difficulty{1}
For every natural number $\ell$, the work assigned to the binary factor search is
$2^{2\ell}$ times the sum of the binary product comparison work at length $\ell$ and
$4$.
\end{theorem}

\begin{theorem}[Modulus Search Work Formula]
\lean{FABL.binaryModulusSearchWork_eq}
\label{FABL.binaryModulusSearchWork_eq}
\leanok
\uses{FABL.binaryConstructorTraversalWork_eq, FABL.binaryFactorSearchWork, FABL.binaryModulusSearchWork, FABL.length_allBinaryVectors}
\difficulty{1}
For every natural number $\ell$, the binary modulus search work equals $2^\ell$ times
one more than the binary factor search work.
\end{theorem}

\begin{theorem}[Operation table work formula]
\lean{FABL.binaryOperationTableWork_eq}
\label{FABL.binaryOperationTableWork_eq}
\leanok
\uses{FABL.binaryConstructorTraversalWork_eq, FABL.binaryOperationTableWork, FABL.length_binaryVectorPairs}
\difficulty{1}
For every dimension $\ell$ and every operation with work cost $w$, the work to construct
its table on pairs of binary vectors of length $\ell$ is $2^{2\ell}(w+1)$.
\end{theorem}

\begin{theorem}[Binary addition table work]
\lean{FABL.binaryAdditionTableWork_eq}
\label{FABL.binaryAdditionTableWork_eq}
\leanok
\uses{FABL.binaryAddWork_eq, FABL.binaryAdditionTableWork, FABL.binaryOperationTableWork_eq}
\difficulty{1}
For every natural number $\ell$, the work assigned to the full binary addition table on
vectors of length $\ell$ is $2^{2\ell}(\ell+1)$.
\end{theorem}

\begin{theorem}[Work of the multiplication table]
\lean{FABL.binaryMultiplicationTableWork_eq}
\label{FABL.binaryMultiplicationTableWork_eq}
\leanok
\uses{FABL.binaryMulModWork_eq, FABL.binaryMultiplicationTableWork, FABL.binaryOperationTableWork_eq}
\difficulty{1}
For every natural number $\ell$, the work assigned to the multiplication table on binary
vectors of length $\ell$ is $2^{2\ell}\bigl(\ell(6\ell+2)+1\bigr)$.
\end{theorem}

\begin{theorem}[Binary field representation bit count]
\lean{FABL.binaryFieldRepresentationBits_eq}
\label{FABL.binaryFieldRepresentationBits_eq}
\leanok
\uses{FABL.binaryFieldRepresentationBits, FABL.length_allBinaryVectors, FABL.length_binaryVectorPairs}
\difficulty{3}
For every natural number $\ell$, the number of bits in the binary field representation
equals $\ell\bigl(1+2^\ell+6\cdot 2^{2\ell}\bigr)$.
\end{theorem}

\begin{theorem}[Successor bounded by a power of two]
\lean{FABL.nat_succ_le_two_pow_succ}
\label{FABL.nat_succ_le_two_pow_succ}
\leanok
\difficulty{4}
For every natural number $\ell$, one has $\ell+1\le2^{\ell+1}$.
\end{theorem}

\begin{theorem}[Product comparison work bound]
\lean{FABL.binaryProductComparisonWork_le}
\label{FABL.binaryProductComparisonWork_le}
\leanok
\uses{FABL.binaryProductComparisonWork, FABL.binaryProductComparisonWork_eq}
\difficulty{3}
For every natural number $\ell$, the product comparison work is at most $4(\ell+1)^2$.
\end{theorem}

\begin{theorem}[Binary factor search work bound]
\lean{FABL.binaryFactorSearchWork_le}
\label{FABL.binaryFactorSearchWork_le}
\leanok
\uses{FABL.binaryFactorSearchWork, FABL.binaryFactorSearchWork_eq, FABL.binaryProductComparisonWork, FABL.binaryProductComparisonWork_le}
\difficulty{4}
For every natural number $\ell$, the work required by the binary factor search is at
most $8(\ell+1)^2 2^{2\ell}$.
\end{theorem}

\begin{theorem}[Modulus Search Work Bound]
\lean{FABL.binaryModulusSearchWork_le}
\label{FABL.binaryModulusSearchWork_le}
\leanok
\uses{FABL.binaryFactorSearchWork, FABL.binaryFactorSearchWork_le, FABL.binaryModulusSearchWork, FABL.binaryModulusSearchWork_eq}
\difficulty{5}
For every natural number $\ell$, the binary modulus search work is at most $9(\ell+1)^2
2^{3\ell}$.
\end{theorem}

\begin{theorem}[Bound on binary addition table work]
\lean{FABL.binaryAdditionTableWork_le}
\label{FABL.binaryAdditionTableWork_le}
\leanok
\uses{FABL.binaryAdditionTableWork, FABL.binaryAdditionTableWork_eq}
\difficulty{4}
For every natural number $\ell$, the work assigned to the full binary addition table on
vectors of length $\ell$ is at most $(\ell+1)^2 2^{2\ell}$.
\end{theorem}

\begin{theorem}[Multiplication table work bound]
\lean{FABL.binaryMultiplicationTableWork_le}
\label{FABL.binaryMultiplicationTableWork_le}
\leanok
\uses{FABL.binaryAddWork, FABL.binaryArithmeticWork_le, FABL.binaryMulModWork, FABL.binaryMulModWork_eq, FABL.binaryMultiplicationTableWork, FABL.binaryMultiplicationTableWork_eq}
\difficulty{4}
For every natural number $\ell$, the work required to construct the binary
multiplication table of dimension $\ell$ is at most $9(\ell+1)^2 2^{2\ell}$.
\end{theorem}

\begin{theorem}[Binary field representation bit bound]
\lean{FABL.binaryFieldRepresentationBits_le}
\label{FABL.binaryFieldRepresentationBits_le}
\leanok
\uses{FABL.binaryFieldRepresentationBits, FABL.binaryFieldRepresentationBits_eq}
\difficulty{4}
For every natural number $\ell$, the number of bits in the binary field representation
is at most $8(\ell+1)^2 2^{2\ell}$.
\end{theorem}

\begin{theorem}[Representation bits within preprocessing work]
\lean{FABL.binaryFieldRepresentationBits_le_preprocessingWork}
\label{FABL.binaryFieldRepresentationBits_le_preprocessingWork}
\leanok
\uses{FABL.binaryFieldPreprocessingWork, FABL.binaryFieldRepresentationBits}
\difficulty{3}
For every natural number $\ell$, the number of bits in the binary field representation
is at most the total binary field preprocessing work.
\end{theorem}

\begin{theorem}[Binary field preprocessing work bound]
\lean{FABL.binaryFieldPreprocessingWork_le}
\label{FABL.binaryFieldPreprocessingWork_le}
\leanok
\uses{FABL.binaryAdditionTableWork, FABL.binaryAdditionTableWork_le, FABL.binaryFieldPreprocessingWork, FABL.binaryFieldRepresentationBits, FABL.binaryFieldRepresentationBits_le, FABL.binaryModulusSearchWork, FABL.binaryModulusSearchWork_le, FABL.binaryMultiplicationTableWork, FABL.binaryMultiplicationTableWork_le, FABL.binaryVectorEnumerationWork, FABL.binaryVectorEnumerationWork_le, FABL.nat_succ_le_two_pow_succ}
\difficulty{5}
For every natural number $\ell$, the binary field preprocessing work is at most
$2^{8(\ell+1)}$.
\end{theorem}

\begin{theorem}[Quadratic bound on arithmetic work]
\lean{FABL.binaryArithmeticWork_isBigO}
\label{FABL.binaryArithmeticWork_isBigO}
\leanok
\uses{FABL.binaryArithmeticWork, FABL.binaryArithmeticWork_le}
\difficulty{3}
As the natural number $\ell$ tends to infinity, the binary arithmetic work is $O((\ell+1)^2)$.
\end{theorem}

\begin{theorem}[Asymptotic preprocessing work bound]
\lean{FABL.binaryFieldPreprocessingWork_isBigO}
\label{FABL.binaryFieldPreprocessingWork_isBigO}
\leanok
\uses{FABL.binaryFieldPreprocessingWork, FABL.binaryFieldPreprocessingWork_le}
\difficulty{3}
As $\ell$ tends to infinity through the natural numbers, the binary field preprocessing
work is $O(2^{8\ell})$.
\end{theorem}

\begin{theorem}[Binary field model resource bounds]
\lean{FABL.buildExecutableBinaryFieldModel_resource_bounds}
\label{FABL.buildExecutableBinaryFieldModel_resource_bounds}
\leanok
\uses{FABL.binaryArithmeticWork, FABL.binaryArithmeticWork_le, FABL.binaryFieldPreprocessingWork, FABL.binaryFieldPreprocessingWork_le, FABL.binaryFieldRepresentationBits, FABL.binaryFieldRepresentationBits_le_preprocessingWork, FABL.buildCertifiedBinaryFieldImplementation, FABL.buildExecutableBinaryFieldModel, FABL.length_allBinaryVectors, FABL.length_binaryAdditionTable, FABL.length_binaryMultiplicationTable}
\difficulty{3}
For every positive integer $\ell$, the constructed model has $2^\ell$
elements and each of its addition and multiplication tables has $2^{2\ell}$
entries. The number of bits in its field representation is at most the preprocessing
work, which is at most $2^{8(\ell+1)}$, and its arithmetic work is at most
$8(\ell+1)^2$.
\end{theorem}
